
% 来源：EPUB part227.xhtml
\section{2025年浙江大学问答实录}

（2025年1月5日）

有机会给学校做点贡献，我自己也很开心。要说分享，我觉得这三条是最重要的。

第一，做对的事情，把事情做对。很多人可能会问，我不知道什么是对的事情。总会知道的！你知道什么是错的，当你发现错了的时候，要及时改。不管付出了多大代价，都往往是最小的代价。

第二，胸无“大”志。那个“大”字是打了引号的，我觉得就是要脚踏实地地去做事情。很多人说，我要创业，创业的目的就是为了要创业。如果你心里没有一个你真的想做的事情，只是为了创业而创业的话，成功率会非常低。

第三，要做一个正直的人。就是要有是非心，要正直，要对得起自己的未来。那么这有没有好处？我不知道。但是，你自己会觉得比较坦然。

% 来源：EPUB part228.xhtml
\BookMinorHeading{如何顺应AI时代}

\BookQuestion{主持人：AI背景下，我们学生在平常的学习中以及未来的人生规划上，可以有哪些顺应时代的想法？}

其实时代一直都在变，工具在变，但是基本的东西还是一样的。大家在学校是学会学习的方法。有AI只是帮你提高效率。当然做论文的时候，大家要小心，不要出“那种事情”。

最主要的，说来说去，可能今天你的好多问题，我都是一个回答：你要想长远，你要想5年以后、10年以后、20年以后，你每一个决策是不是对。学习方法也好，创业也好，我觉得所有方面其实都是相通的。

\BookQuestion{主持人：感谢段学长，不仅给了我们一些建议，也警示了要适当使用AI工具。}

运用AI工具，尤其是做学术、做研究，大家还是要谨慎。

% 来源：EPUB part229.xhtml
\BookMinorHeading{投资做得好的人都很慢}

\BookQuestion{主持人：在投资时，如何快速判断一家公司是否值得长期投资？}

基本上我不太会这个事情，你看我过去这十多年，说来说去就那几家公司。那也是我很多年对企业、对生意、对产品，对所有东西的理解攒下来的。我也没见过谁真的能够那么快，包括巴菲特也好、芒格也好，其实投资做得好的人都很慢。大家并不在乎失去一些机会，但是，最重要的是你不要去踩雷。

很多人也可能会说，你当然不需要挣快钱，因为你有钱。我们需要挣快钱，因为我们没钱啊。这很可能就是你没钱的原因，因为老想挣快钱。其实我也想挣快钱，谁不想挣快钱？但是知其不可为很重要，就是你还是要踏踏实实做该做的。

\BookQuestion{主持人：那段学长的意思就是要敢于尝试，并且眼光要放得远一些。}

哎，敢于尝试这个……要看你怎么个尝试法。我不是一个不敢冒风险的人，但是要冒能够承受得起的风险。大家可能会说，谁谁谁他就赌对了。那我最厉害的是20分钟赚了20多倍。人家说你干嘛呢？我在赌场用100块钱赚到2000多块钱就走了。但是我可以重复这件事吗？不能。

我当年投网易，也是几个月赚了二十几倍，真是二十几倍，而且我是全仓的。人家说你真厉害，你再来一次？我说不会，我是碰上的。苹果我们投得也很好啊，我是2011年投的，现在都2025年了，差不多刚好14年，看起来是很好。但你要是对生意理解不了，对文化理解不了，对商业模式理解不了，早就跑了，也不可能留到现在。所以快这个东西是比较难的。竞技比赛的运动员是需要快的。

% 来源：EPUB part230.xhtml
\BookMinorHeading{该创业的人根本不需要鼓励}

\BookQuestion{主持人：作为一名创业者与投资者，您还会鼓励当代的大学生创业吗？为什么？}

毕业前还是毕业后？

\BookQuestion{主持人：毕业前在学校里有一些创意创业的想法，在毕业后具体去落地。其实我们也有很多同学在学校中开展一些创业项目。}

我没有说过，就不叫“还会”啊，我没有鼓励过。我觉得该创业的人根本就不需要你鼓励。你有个很强烈的想法，就一定会去做。

现在的创业条件比我们当年要好很多了，还能找到VC、Startup这种各种各样基金的支持。我们那个时候都是从零开始的，没有什么别的东西可以帮你，只能靠自己一点一点做。所以现在和那个时候确实是不完全一样，现在创业条件要比那个时候好。但是很重要的一点就是大家不能为了创业而创业，你得是真的有想法。还有走投无路的时候，也是创业的好时机，我们当年就是走投无路。

% 来源：EPUB part231.xhtml
\BookMinorHeading{学习“学习的能力”}

\BookQuestion{主持人：请问段学长，在不断变化的时代背景下，您个人是如何保持学习和适应新知的呢？}

非要学新东西吗？是这样的，在学校，你主要还是学习“学习的能力”。当你碰到有兴趣的方向，你就不会对未来感到恐惧。你觉得有兴趣，就可以想起来我该怎么学习。你会去找书、查资料，你有理解力，包括逻辑能力，这种学习能力与受教育程度低的人相比，相对来讲要强很多。

比方说我妈妈，iPhone不肯用，iPad也不肯用，她说这个太难了，我肯定不会。她是拒绝学习的，好多人是拒绝学习的。前段时间我看了我一同事写的文档，我说你这个是手写板写的吧？他说是。我说你为什么不练打字呢？哎呀，他说，不一直都没练嘛。他比我小10岁，我说你这个没有道理，我花了3天的时间就学会打字了。我也用过写字板，接在电脑上，太不方便了。其实你只要不恐惧，知道自己能学会，就没有什么问题。当然，你还是要学你需要的东西，这是一个知识爆炸的时代，不能说所有新东西你都学，这不得累死，学了也白学。很多人看起来知识很渊博，你再渊博，还能渊博得过搜索吗？所以，何苦呢。

% 来源：EPUB part232.xhtml
\BookMinorHeading{做企业、做投资都是游戏}

\BookQuestion{主持人：一个好赛道往往会随着竞争的加剧而进入低毛利的时期。那么在低毛利时期，创业者该如何应对呢？}

首先，好赛道是不会进入低毛利时期的，低毛利的都是商业模式比较差，产品差异化很小的。作为创业者，你还要进去，那就脑子坏了呗。

投资也好，创业也好，其实都是要很认真地去想商业模式的。但我觉得创业还是有点不一样，你必须要创业的时候，你要做的东西是你自己真的有感觉的。我当年做游戏没有别的原因，就因为我自己爱玩游戏。所以我从来不批评爱玩游戏的学生，我就那样，我儿子要玩游戏，我让他玩。我只是想办法跟他商量，你玩多少时间，就要干很多事去挣时间，而不是说简单地不让玩。我觉得所有的事其实都是游戏。有人喜欢学习，其实是他有快乐嘛，游戏也是可以带来快乐的东西。

反正创业是这样的，你要是完全想赚钱，想出人头地，但不知道自己能做什么，那真不知道最后能做出啥来。我做游戏的时候，像做小霸王游戏机时，我知道很多人会像我一样爱玩游戏。那我想的事情就特简单，我并不需要说服别人去玩游戏，我只要把产品做好就好。你要搞一个新东西，要让别人知道这是个啥，那个过程远比你把大家都喜欢的东西做好——质量做好了、渠道做好了、服务做好了——难度要大得多。后面我们做企业，其实没有那么难。当然了，当规模大到一定程度，你就必须得往前走。我们早年说敢为天下后，那时候我们很小，现在我们也做很多新东西。没办法，你不做，你前面就没人了，那怎么办？

% 来源：EPUB part233.xhtml
\BookMinorHeading{做对的事情就是发现错了马上改}

\BookQuestion{主持人：您的本分哲学中说到做对的事情，那么大学生要如何判断一件事是否是对的呢？}

其实是知道的。有些事情你不知道是否对，那你做了，你将来会知道是否对，因为如果不对，你会受到惩罚。但是呢，你发现错了，一定要马上停。大部分事情，我想大部分人是知道对错的，但还是有很多人做错的事情。为什么很多人明知是错，他还要继续做呢？因为错的事情往往有短期的诱惑。比方说抽烟、酗酒等等就是这样，你看多少人该抽还抽，难道他们不知道那个不好吗？他们知道。当然这样的事情不伤害别人，也不是问题，只是说对自己身体不好。还有比方说舞弊的、论文作假的，他们不知道这是错的吗？他们知道。但为什么还要做？难道这个也要大学教？我还见过以前有些农村学校写的标语，说“不许强奸学生”，这怎么能写得出来呢？他们当然是知道对错的，然后非要靠法律来处置这样的事情，那就很难办了。所以，为什么说要正直，就是你自己要有一个原则去约束自己，你发现错了，就马上改。

生意上也是一样，有人买了只错的股票，如果他是炒股票我就不说了。如果他原来以为这个公司是这样的，结果发现是那样的，他不马上跑，他还说，等它涨回来我再跑。我记得有个人跟我讲乐视，说乐视现在掉了很多，到三十几块了，问我怎么办。我说三十几块的价钱很好啊，因为将来会是零的嘛，所以你现在什么价格卖都是好价钱，你要等它涨回五十块，这个有点难。当然我后来发现，他是想把这股票卖给我，我觉得这是脑子有病。他是原来乐视的第二大股东，我觉得那是个坏人。

% 来源：EPUB part234.xhtml
\BookMinorHeading{老师真的关心学生，关系自然会好一些}

\BookQuestion{主持人：您求学期间与信电老师深厚的师生情谊令人动容，请问您认为一名刚迈入职场的青年教师，该如何传承好这份育人精神，实现与学生的双向奔赴？}

这个有点难，我对现在的学校生活已经不太了解了，但我觉得老师跟学生，关键应该是在老师吧。老师真的关心学生，那跟学生的关系自然就会好一些。我搞不清楚现在老师跟学生有多少接触，我们那时候确实是多一些的。学校也少，人也少一些嘛，地方也小。我们就在三分部，后来是大一的学生去的是吧？现在呢，现在都不去了？

\BookQuestion{主持人：现在是法学院。}

% 来源：EPUB part235.xhtml
\BookMinorHeading{相信中美关系将来也会好起来}

\BookQuestion{主持人：请问段学长在中美对立的时代，浙大人的创业机会在哪？作为校友，能够为创业的校友提供哪些支持？}

中美对立吗？我们学校谁建的？大家不要忘了我们读的浙江大学，是有很多老美的痕迹的。大家对这个东西不应该忘记，其实人家对我们还挺好的。我看很多人现在去美国，没有太多的障碍。我相信将来也会好起来。可能眼前有一些短期的冲突，包括又选了这么个总统，这确实是会有一些麻烦。

% 来源：EPUB part236.xhtml
\BookMinorHeading{创业成功率非常低}

\BookQuestion{主持人：如果您今天是20岁，要创业，您会选择哪个赛道以及选择合伙人的标准是什么？}

这个还真不知道，我觉得如果我有机会，我还是找份工作，好好享受人生。创业并不容易的，不要看到有些人成功，成功率非常低。大部分人忙一辈子，其实没什么好的结果。你看竺可桢学院本身也很厉害，黄峥就从那出来的，我知道他们厉害，但是也不意味着他们成功率就很高。不过想去创业，你就尽管去，现在条件真是比我们原来那个时候好很多。

\BookQuestion{主持人：谢谢段学长的回答，我猜测刚才这位同学应该也是有创业的想法。所以想要来直接提问一下。}

有想法有条件就去做，你如果只是为了创业而创业，你就小心，这就是我的观点。我并不知道大家应该做什么，现在哪知道啊，信息这么爆炸，我整天忙的就是打球，谁还管这些事。

% 来源：EPUB part237.xhtml
\BookMinorHeading{不清楚现在创业从哪搞}

\BookQuestion{主持人：请问段学长，文科学生创新创业可以走怎样的赛道？这是我的一个问题，因为我作为一个法学专业的同学，本科期间接触创新创业是比较少的。}

简单讲，我不知道，因为我不是读文科的，我连语文课都没上过，真没上过。我们中学没有语文课，只有政治语文课，学的都是那些什么“天连五岭银锄落，地动三河铁臂摇”（毛泽东《七律二首·送瘟神》）。你们太年轻了，不一定听得懂，所以是不一样的。我上大学的时候，语文居然考了60多分，我还蛮意外的，确实是没有任何底子。但是我妈妈在图书馆工作的时候，我老是去图书馆，虽然看的书没有那么高级，也过得去。《金光大道》《艳阳天》，你们听起来估计都很陌生，你们太年轻了。所以我不清楚现在创业从哪搞，真的搞不清楚。

% 来源：EPUB part238.xhtml
\BookMinorHeading{养成安全运动习惯}

\BookQuestion{主持人：请问段学长平常累的时候会怎么放松？}

累的时候就该睡觉了吧？睡得好的话，你可以多拉伸，泡个热水澡，都是属于物理治疗，没有什么特别的。我们都是凡人，都是普通人，我这里没有什么神秘的东西，没有什么诀窍。

\BookQuestion{主持人：我们的西区马上也要建成了，是由段学长冠名的。段永平文体中心是一个非常大的文体活动中心，在未来会有非常多的同学在这个文体活动中心活动，包括各类锻炼和放松，也欢迎未来段永平学长多回来，多回母校看看。}

今天我去看了那个地方，我觉得挺好的。同学们在上大学时，当然包括大学以前，要养成运动的习惯，安全地运动，要学会拉伸。我们以前是不知道的，搞得我现在有时都觉得膝盖有点小问题了，以前不够注重这些。现在条件好，而且里面有人教，挺好的。

\BookQuestion{主持人：我们浙大同学的课余文体活动非常丰富，非常欢迎段学长到时候到建好的场馆里面，和同学们一起打打球之类的。}

我还能打点乒乓球。

% 来源：EPUB part239.xhtml
\BookMinorHeading{最重要的是做对的事情，把事情做对}

\BookQuestion{主持人：请问段学长，您觉得浙大有什么样的气质和氛围影响了您的创业、投资和生活呢？}

我记得我们单位曾经还有一个人说，一定要跟着段永平的脚步，所以要把儿子送到信电系。我说那我下过哪条河摸鱼，你没有去过，你也不知道。影响成长的因素有很多种，我的成长跟浙大有没有关系，当然是有的。浙大毫无疑问是个好学校，但是也会有问题。我对现在不了解，现在学生和老师都多了好多。我在三分部的时候，师生关系也挺好的。

\BookQuestion{主持人：现在的信电学院，也是段学长您的母院，有特别经典的两个词，叫作勤奋乐观。我们要争做勤奋乐观的信电人，这两个词也是信电学院整体学生的特质。在学习上，我们要时刻保持勤奋，保持不畏难，克服更多的难点。在平常的生活中，我们要保持乐观向上的态度。}

勤奋和乐观？ 我觉得最重要的是要做对的事情， 把事情做对。

我要你勤奋， 你勤奋得起来吗？ 我就从来不勤奋。我可是想尽办法偷懒， 我偷懒的意思不是说， 考试偷看、用ChatGPT来写论文， 你参考是可以的。我觉得最主要还是脑子里要有是非， 要想着做对的事情， 发现错了， 要赶紧改。乐观？ 乐观是个性的问题吧， 不是每个人都能想乐观就乐观得起来的。这两个词说完了以后， 大家真的会得到什么好处吗？ 我觉得可以质疑一下。

% 来源：EPUB part240.xhtml
\BookMinorHeading{创新都是从模仿开始的}

\BookQuestion{主持人： 当然在勤奋乐观之前， 更多的先是我们浙大的求是、创新的精神。刚才老师送您的求是卷轴……}

创新也是要小心的。很多人把创新理解成不一样， 我觉得那个要特别小心。最主要是你到底想干什么， 你给用户带来什么价值。你说我做了一个东西， 以前没人做过的， 创新了吧？ 可是也没有人喜欢， 你不就惨了吗？ 所以， 创新是你满足了人家的需要， 也是别人没有做过的事情。

差异化不是说不同， 差异化其实跟创新是一个意思， 就是你要去满足用户的需求。用户包括很多方面， 老师满足学生的需求， 学生也是用户。例子很多了， 建楼质量建好了， 那也是一种服务， 你如果只是想建得好看， 但是不好用， 不就瞎干了嘛？ 扯远了。

\BookQuestion{主持人： "敢为天下后， 后中争先" 这其实也是学长的理念， 就是不能去盲目创新， 是吧？}

我说过敢为天下后， 我可以给你举很多例子。比方说谷歌敢为天下后、苹果敢为天下后、微软敢为天下后， 那你说还有谁不是？ 做得成功的企业早期都是这样。这些年可能有些新的， 比方说AI， 其实很多人也是跟着做的。你看豆包， 它也不是新的。多少人做chat是从ChatGPT开始。但是这个其实我也搞不清楚， 我现在用好多个工具， 像ChatGPT、Gemini、Perplexity， 豆包也用， 你说他们创新吗？ 创， 但是他们其实未必是最先做的。谁发现了这个东西，又刚好有这个实力，把它做得比别人更好，他可能就厉害了。搜索显然不是谷歌先做的，雅虎肯定是早就做了。苹果的iPhone，我们做手机感觉都比它早呀，可能真的比它早，但人家敢为天下后，它实力强，后发制人了。

但是有些东西呢，后也不一定行。哪怕是微软，它的搜索就没有办法搞过谷歌，所以它就特别热衷于搞ChatGPT，因为想从这个角度把搜索抢回来。这是我瞎猜的，我不知道，我也没那么关心，我就想着好好打打球啥的。

\BookQuestion{主持人：请问您如何看待追热点，或者说模仿与创新之间的关系？}

追热点，这个定义不是很完整，我也不知道什么叫追热点。追星也算，追风口也算。你刚刚说模仿与创新，你要能够做出不模仿又是很好的东西，那当然可以。但是我几乎都很难找出一个例子，它不是从模仿开始的。我刚刚讲的微软也是，微软就没有什么东西是一开始就做的，都是有类似的产品在前面，做完了以后它跟着做的。苹果也是，谷歌也是，亚马逊其实也是。脸书算不算我还不好说，因为那个时候社交媒体好像没有那么多，但它其实也是在别人的基础上建立起来的，好像是在哈佛的校园网上。所以我觉得，这不是个矛盾，没有什么冲突在里头。

% 来源：EPUB part241.xhtml
\BookMinorHeading{蚊子还是那些蚊子}

\BookQuestion{主持人：请问段学长，您记忆中的校园和今天看到的有哪些相同和不同？}

说句实话，我也没怎么注意。我坐车晃了一圈，也没有下去散步。应该没有什么太大的变化，蚊子还是那些蚊子。

% 来源：EPUB part242.xhtml
\BookMinorHeading{养成想长远、想本质的习惯}

\BookQuestion{主持人：请问段学长，如何培养看待问题的批判性思维？}

这个我搞不清楚，我没有什么批判性思维。你想长远、想本质嘛。

你跟我做得一样，你挺好，我干嘛要批判你？但是有些东西，像你刚才讲的勤奋乐观，确实有人说过勤奋最重要，我说做对的事情最重要。因为你说勤奋最重要，你会变得更勤奋吗？你不会，至少我不会。

但是你选择做对的事情，发现错了，你赶紧改。这一辈子的每一个决定都是基于10年、20年以后回过头来看，你会发现你会省好多力气。我们毕业三四十年了，我看那些同学都是很聪明的人，但有些人在原地打转了30年、40年。可能做每个小的决定，他都是基于眼前的利益，他们可勤奋了，但是不管用。我有时候会想这个事，所以养成了想本质的习惯，我觉得挺好，我不觉得是批判性思维。凡事就这么想，时间长了，你可能就有这种能够想本质的习惯。但是不是叫批判性的思维我就不知道了，它听起来好像恶狠狠的样子。

% 来源：EPUB part243.xhtml
\BookMinorHeading{对抗压力的办法}

\BookQuestion{主持人：如何成为像您这样一个非常坚定，特别高能量的人，能够欣然应对或者去承受工作中、生活中的挑战和痛苦？}

我是凡人，我也一样有痛苦。而我还是个懒人，你看我早早就退休了。我退休20多年了，我今年63岁，大概40岁就退休了。显然，我就是不愿意面对压力。我已经享受过了，我知道工作的乐趣。然后呢，我们有更年轻、更能干的人，他比我干得还好，何乐而不为呢？但这不是我想教你们的，不是每个人都有那样的运气，正好有那么多好伙伴可以跟你一起合作。

当然，我也有我的心得。我能做到今天，是因为我能支持别人，我有分享的精神。很多人在用人的时候，都是要让别人为自己工作，我是希望他们能够为他们自己，而不是为我工作。包括，要花很多时间去建立一个系统，这样我就可以面临小一点的压力，而不是去硬扛。

举个特别简单的例子，我们公司实际上是没有销售部的，我们的产品只是给代理。因为早年我发现，销售部的工作就是每个客人都来找你谈呀谈呀，谈来谈去全是谈价钱。后来我就想，必须得解决这个问题。所以我们公司的产品价钱对每个人都是一样，没有说客户大，就有折扣，那得多累啊。这个压力很大，我那时一天可能吃8顿饭，去6次桑拿，去五六七八次的卡拉OK，神经病那是。我肯定不想去，但没办法，每个客户都缠着你，要跟你谈价钱。而且那个时候我们才做很小的生意，就已经是那样了。我就想10年、20年以后你不就完蛋了吗？当然要那样忙的话，也可能熬不到10年20年。我们花了大概3年的时间，建立了销售系统，这也是一种对抗压力的办法。

% 来源：EPUB part244.xhtml
\BookMinorHeading{回到事物的本源}

\BookQuestion{主持人：刚才您也提到了，年轻人要为自己而工作。那么顺着您的回答，刚好下一个问题：年轻人的什么特质会吸引到您，以及您希望会与怎么样的人成为合伙人？}

你说我认为年轻人要为自己工作，我不是那个意思。我的意思是，我要让我们的员工为他们工作，这个意思上好像是有点差别。当然实现自我也是应该的，但是跟我说那个话的意思不太一样。因为作为一个CEO，我要去支持他们。我让他们觉得他们是在为自己工作，而不是跟他们说，你们要为自己工作，然后就把我给踢走了，意思不完全一样。

\BookQuestion{主持人：年轻人的什么特质会吸引到您？}

你是指哪方面的？年轻人年轻啊，你们的年龄就是我的梦想，现在我都回不去了啊。

\BookQuestion{听众：您和黄峥之间有很深的友谊，你可能看重黄峥身上的一些点。所以我想请问您，年轻人的哪些特质会非常吸引您？}

其实跟年纪没有关系，我肯定不喜欢跟我没有办法沟通的人打交道。我非常喜欢朋友，跟朋友聊天我很开心。我觉得黄峥是一个看本质的人，跟我一样，我觉得这还是比较重要的。他到底算不算年轻人？他当然算，其实我也算。又回到我前面说的：要去做对的事情，要把事情做对，要想长远。我觉得黄峥他就是这么一个人，所以我们在一块聊天的时候，我们不会聊着聊着就不在一个频道上了。但大部分人是会的，我跟很多人聊天，其实都不容易聊得下去。不知道算不算回答了你这个问题，我觉得还是要回到事物的本源上。

% 来源：EPUB part245.xhtml
\BookMinorHeading{坚持做对的事情}

\BookQuestion{主持人：假如您刚刚大学毕业，选择一件事情每天坚持做，您会选择什么？}

睡觉。每天坚持是啥意思啊？经常做，比如体育锻炼要做。还有啥？太多事情需要做了，但是必须要选一件事情做，或者两件事情，这个问题有点难。

\BookQuestion{主持人：其实是想问您习惯要如何养成。}

运动很重要，也不一定天天了，因为你要保持身体的健康、保持大脑的清醒。然后又回到我刚刚说的，逻辑上要比较清楚，凡事要想长远，要坚持做对的事情，发现错了，要赶紧改。把事情做对是个学习的过程，就是说你会犯错误。有很多人分不清楚错的事情和把事情做错，这两个性质是完全不一样的。不能因为怕把事情做错了就不做事。因为无论你怎样做，都可能会错。但明知是错的，你不应该做。

% 来源：EPUB part246.xhtml
\BookMinorHeading{不要把重要变成紧急}

\BookQuestion{主持人：可以分享一下您在求学和创业过程中遇到过的挫折吗？}

求学的过程中，我差点大学都没毕业，王东潘老师帮了我很大的忙，不然我连答辩都过不去，真有可能毕不了业，那也算挫折了。

创业过程中的挫折，那太多了，每一个产品的问题都是。我那时候的电话是24小时开机，但是我非常不愿意接电话，因为每一个电话可能都是麻烦，太多了。

\BookQuestion{主持人：面对这些挫折的时候，有什么解决的心得，可以跟我们分享吗？}

其实还是有的，做事情你要尽量地想长远，又回去了。你不要让重要的事情变成紧急的事情，我觉得这个很重要。一天到晚接电话这个问题，我现在就没有了，我记得很久以前，我就基本上不用接电话。要让自己不要有太多紧急的事情。比方说要锻炼身体，比方说不能等到病得不行了才去看医生。看医生也很重要，定期的检查很重要。

我以前在北京做过一个节目，叫《危机时刻》，我不知道你们有没有注意。那时候我当CEO，其实经常接受采访，为了企业嘛。主持人问我，危机时刻你怎么办？我说我开着车，时速200千米，前面20米有一堵墙，肯定要撞上去了，你觉得我该怎么办？没什么办法，反正死定了。最重要的是，你不要开那么快呀，你干嘛要去撞墙呢？你开辆坦克才行啊。意思就是要完全从长远的角度去解决这个问题，不能等到危机时刻才想办法。大家很喜欢看《危机时刻》，觉得哇，这个人临危不惧。最后不还是死了嘛，并没有意义啊。你要预防这东西，所以安全第一。安全第一并不是说到危机时刻你有办法，那个时候没办法的，你要提早做准备。

% 来源：EPUB part247.xhtml
\BookMinorHeading{除了价值投资，没有其他投资方法}

\BookQuestion{主持人：请问在您的视角中，信息差对于投资和选择的影响有多大？}

炒股票的，对我没有什么影响，除非你买了马上想卖，想赚点那种不该赚的钱。当然了，也有一些是靠信息差赚钱的，我没有那么喜欢这些事情。对我来讲，投资不是一个零和游戏。信息差实际上是个零和游戏，就是利用这个信息，在别人没得到之前，你赚了别人该赚的钱或者怎么样。量化投资是有点那个味道，但是你不做，别人做，它作为一个生意无可厚非。

但你看长远以后，这些东西都是小小的波澜。你太在乎了，会很累，球还要不要打了？要生活的嘛。所以要看长远，找到好公司你就拿着。那“拿着”听见了，“好公司”就没听见。拿个错公司，不就死得更惨吗？我跟好多人讲，他们就是不明白，他们说价值投资就是长期投资，不是的。除了价值投资，请问还有别的投资办法吗？其实是没有的。你投的就是价值，如果你不投价值，投的是什么？所以看长远，拿到好公司，你的生活可以愉快很多。你看沃伦·巴菲特、查理·芒格，他们活得都很长。你看那些投机很有名的人，大部分都死得很惨，这是真事啊，或者说生意都死得很惨。国内也不是没有啊，一个有钱人，好好地跑去坐了牢了。

% 来源：EPUB part248.xhtml
\BookMinorHeading{不能用需要的钱去赌不需要的钱}

\BookQuestion{主持人：您是否经历过面临高风险与高回报的关键投资决策？面对高度不确定的市场，您如何平衡风险与收益呢？}

明知高风险，还要去做，脑子坏了吧！真是这样的。但是风投不一样，风投他们不冒风险，他是拿别人的钱去投，赚到了钱，有他的一份，他当然可以投。当然他也会有一些判别，看现在或者未来的市场是不是有可能喜欢，他可以融到更多的钱；是第一轮投进去，可以赚第二轮、第三轮、第四轮的钱。风投严格讲是一个很现代的产物，我们那个时候是没有这个东西的。投资本身的逻辑是要看清企业的商业模式、未来现金流，确实能挣到钱你才敢投，这样你也愉快。

但风投不会像我们那样投资，他都是这里放一点，那里放一点，其实赚的是一个国家的钱。你看风投做得很大的，比方说像孙正义，扣掉他前两名赚钱的公司，估计他就赚不到钱了。他其实也没赚多少，但他是非常有名的风投。所以很有意思，他赚的钱跟巴菲特和芒格相比，那确实还是差很多的。很多人做企业，企业做起来以后，作为最早的创始人，在企业市值很高的时候确实身价很高。但那是一个剩下来的人，同一年创业的其实可能有1000万个人。大家都看见马斯克了，也跑去做电动车，但你看看中国一年死多少电动车企业，就是这个道理。

\BookQuestion{听众：举例来说，医药公司的研发面临很高的不确定性，技术能不能做出来，做出来之后又要考虑能不能过临床之类的。这个时候你需要去考虑高风险、高回报的情况。}

你这么说我能够理解。我做企业的时候有可能会去做这件事情。投资人呢，作为风投做这种事情也是有道理的。但是大资金这么做就比较困难。不能用你需要的钱去赌你不需要的钱。我一年赚个10\%、 20\%挺好的。一个可能有500\%的回报，但也可能亏光光的事情，我为什么要去做？所以一般我不会做这样的事情。

但是现在这个机制挺好，要有人去做这个事。你说的生物医药就是个特别典型的地方，其实科技也一样，有很多新的想法。但是没有办法马上就有成熟的生意模式，我是不投的，因为我没有那么多时间，我要打球。很多人说你得关心，我说我不关心，等到它好起来了，我再看也可以，我不是个职业投资人。虽然我管的钱比绝大多数职业投资人可能还多，我一个人管的钱可能相当于一个中型的对冲基金，估计更多，但是我并不会像他们那么忙，我看见好多人就三四百个亿，天天忙。我说你们在忙什么，你看我就买了个苹果，一直拿着就完了嘛，或者说人民币我就买了茅台。我就可以打球了，让他们忙吧，我看见很多投资人都很忙。但是我理解你说的这句话，我非常同意，只是我自己一般会避免。

\BookQuestion{听众：这可以理解为是两种不同的风格，但是您更喜欢找那种高确定性的。}

你说的那个是高风险但不知道有没有高回报的，我要知道它有高回报，我是可以投的，这取决于懂不懂。如果不懂，像我买的公司大家也都觉得是高风险，但我不觉得。所以这取决于你有多懂，要投你懂的东西。我一直说不懂不碰，搞不懂的东西，你不要去碰。另外就是不要用杠杆，不要用股票去抵押。哇，觉得股票要涨了，就借好多钱，结果股票在涨之前掉了一下，你就完了。所以不要碰这种事情，这样就容易很多、你还可以打球。

% 来源：EPUB part249.xhtml
\BookMinorHeading{好公司都过得还不错}

\BookQuestion{主持人：请问段学长，当下年轻人应该如何应对经济下行周期呢？}

其实当下，老年人也面临这个问题，跟年轻人有啥关系？没关系。我真不知道该怎么回答，好自为之呗。确实我们面临这个问题，浙大要好一点吧，你还是可以找到好工作的，你还可以先读个研，实在不行再读个博。如果你刚好看见有一个你喜欢的好公司，有机会加入，也是一个很好的选择。

因为我看到就算是经济不是那么好的情况下，好公司的小日子都过得还不错。腾讯、茅台，还有我们公司，其实都挺好。茅台股价掉不等于公司状况不好，其实好得很，酒还是买不到，我们有直购他们的酒，给的量都很小。

% 来源：EPUB part250.xhtml
\BookMinorHeading{全球化是伪命题}

\BookQuestion{主持人：请问中国企业怎么能够真正地做到全球化和本土化？}

哪个企业？所有的企业吗？现在中国算是做好了全球化的可能就一个拼多多，一个抖音，其他好像都是局部的。但是我觉得所谓的全球化像一个伪命题，你并不需要去追求它，到了该去的时候，自然就去了。你根本就没有那个实力，也没有那个需求的时候，走不出去的。不要说别人，我们公司也是一样。早年我刚到美国的时候，我就觉得：哇，超级碗是一个特别好的做广告的地方，我们要是有合适的产品，一定要到这个地方打广告。但我们一直没有做到，我们一直没有找到合适的产品进入这个市场，今天也没有。

但是黄峥做到了。他有次问我，你觉得我们在超级碗那个广告，效果怎么样？我说我看见了，你自己觉得怎么样？他说效果还可以。我说，还可以，你才打两条？！然后第二年他打了四条还是五条，忘记了。所以他就是一下子抓住了重点。现在我问美国的所有朋友，你们知道Temu（拼多多海外版）吗？他们知道，几乎都知道，我家里工作的很多人都用Temu，阿姨、钟点工。在它上面我要买的东西不多，但是有些东西确实有很多人在买。它最后还是要靠产品的质量啊这些东西，但这个模式还是有点厉害的，它的质量也是一点一点在往前走。

% 来源：EPUB part251.xhtml
\BookMinorHeading{做该做的事}

\BookQuestion{听众：您做的广告都很炫目，但您一直很低调，平时在公众场合为公司代言好像也不多。}

我低调吗？有几个企业家敢在雪球上说话的？

我早就退休了，当CEO的时候我出来很多的，你现在还能找到一些我当年的视频。现在我退休了，不在其位，不谋其政，我要打球啊。所以跟高调低调没有什么关系，其实我很正常，我就是一个普通人。我从来没有刻意地高调，也没有刻意地低调过，我就是做我该做的事情。我不在企业一线了，你找我，我觉得不合适，我不想去跟CEO们抢功劳，让大家觉得这还是我的企业，它不是我的了，是他们的，我只是一个打球的人。

\BookQuestion{主持人：您在雪球上的账号昵称是“大道无形我有型”，不知道段学长为什么取这样一个名字？}

这里面没有什么特别的故事，本来就有“大道无形”这个说法，而且我以前在玩梦幻西游的时候，用的就是大道无形。雪球上好像有人注册过了，那我就“大道无形我有型”，就开了个玩笑呗，表示大道虽然是无形，但我们可以有型的嘛。

% 来源：EPUB part252.xhtml
\BookMinorHeading{巴菲特：他告诉我苹果的生意模式比可口可乐好}

\BookQuestion{主持人：您和巴菲特交流之后，你们在价值投资方面的最大共性和不同是什么？}

我们对投资的理解是一样的。就是我刚讲的，投资其实就是价值投资。你不投价值，那你投的是什么？但投机不一样，人家说投机是投资在机会上，你这样讲也OK。

差异是有的，就是我们对不同生意的了解都是不一样的。很多他懂的生意，我不一定懂。但是我懂的生意，他有时候也未必懂。

举一个简单的例子。巴菲特也买了很多苹果，但他大概是在2016年开始买的，我是2011年买的。2018年的时候我曾经跟他聊过一次苹果，还是专门去跟他聊的。反正我们有一个联络，他不知道听谁说，我估计可能听李录说的，说我有一个博客。他说：“你那博客有没有英文版的，我想看看。”我回答：“沃伦，我的英文你又不是不知道，我哪可能用英文写，这是第一。第二我写的好多东西都是从你那来的。”他说OK。我说：“但是，你要是想聊苹果，我可以来跟你聊。”他说：“哦，任何时候，我太想跟你聊了。”我说好，就去了。他让我任何时候路过奥马哈就去，我就想怎么路过呢？就查了地图，哦，去芝加哥是路过奥马哈的，我就赶紧安排了一个去芝加哥的行程。然后给他发了个电子邮件，说我要路过了，他说那你就来。真是这样的。

然后到那个地方的时候，他坐在球会大堂那等我。我们有三个人，吃了饭，就聊了。他人很好，一开始的时候就想办法活跃气氛，就讲他小时候是喝百事可乐，我说我怎么不知道这事，不都说你是喝可口可乐的吗？他说，那个时候百事可乐的价钱只有可口可乐的一半，所以显然我要喝百事。然后我们开始聊到苹果，就准备吃饭了。

我说，我觉得苹果的生意模式，当然我这里主要讲的是iPhone，是比可口可乐要好的。“哇哦”，他说。我做市场出身的人，我可是知道点在哪，这是他关注的点，当然也是我关注的点。为什么说苹果生意模式好？我说：“你看，你说百事可乐便宜一半的价钱，你就喝百事了。你知道安卓手机和苹果比，价钱差多少？但真正的苹果用户谁在乎？没有人在乎，至少在美国是这样，在中国也是一样的。真的用苹果的人，他就用苹果。”然后这个事就聊完了。他说：“我明白了。”那天晚上我们就聊了好多别的乱七八糟的事，没有再聊过苹果。

后来到2022年，我去他们股东大会，他请我去他们的晚宴，在奥马哈的水族馆。我觉得好有意思，在这种地方设晚宴。我想不会搞错吧，我英文不太好，早点去吧，就提前半个小时去，结果很顺利就找到了，路标什么的都很清楚。进到里面是一个大的厅，像礼堂一样，摆了50张左右吃饭用的圆桌。我进去的时候，里面就一张桌子上坐着两个人，就是沃伦·巴菲特和查理·芒格，我走到他们身边时，对面又走过来一个人，叫比尔·盖茨，然后我们四个人坐下来。巴菲特把我介绍给芒格和盖茨说：“This is Ping。”我的名字叫Ping。巴菲特说：“他告诉我，苹果的生意模式比可口可乐好，所以我买了很多苹果。”你想隔了四年，就是我们那一天晚上说的那么一句话，他就抓住重点了。我跟无穷多的人讲过，没有几个人能够抓住的。我跟他平时也不会见，四年以后，他直接就这么介绍了。当然他买了很多苹果，这个跟我没关系，那只是玩笑，但是他拿到现在还没卖完，可能跟我是有关的。当然我没有办法验证这个事情。

% 来源：EPUB part253.xhtml
\BookMinorHeading{在浙大学会了学习的方法}

\BookQuestion{主持人：请问在成为优秀的投资人、公益者的道路上，浙大对您产生了怎样的影响？}

我在这里上了大学啊，我在这里学会了学习的方法，还是有影响的。浙大的学风也应该还可以，主要是它不在大城市里头，所以可玩的地方比较少一点。杭州现在其实已经很厉害了，我们两个公司都在杭州设立分公司，这边的人才很多啊。而且城市到现在为止，人口都是净流入的，很多人才会流进来。所以它很好，但是不是就变得更多地方可以玩了？应该也是更好玩了，这不是什么坏事，我自己也喜欢打游戏。我们那个时候没别的可以玩，我读书的时候也没有游戏，还真没有，顶多下个棋。

% 来源：EPUB part254.xhtml
\BookMinorHeading{好好运动，好好生活}

\BookQuestion{主持人：您的下一个人生目标或者说终极目标是什么？}

健康地活长一点。生活还是要有质量，如果只是活着，插着管儿，为了某个理由，我觉得太痛苦了。所以我理解琼瑶，我觉得她是有道理的，所以该给我拔管的时候就拔吧。但是你要为这个努力啊，你要好好地运动，好好地生活，心情要好。这个时候需要乐观。

% 来源：EPUB part255.xhtml
\BookMinorHeading{对AI不明觉厉}

\BookQuestion{听众：您怎么看待AI的发展？您觉得AI在这几年会不会发展得非常迅猛，又会给我们带来什么样的影响？}

AI发展非常迅速。

我的简单说法叫不明觉厉，我确实觉得它厉害，但确实不是很懂。以后我还要问你们呀，这个AI到底会怎么样，学院的人肯定比我们更靠前，看得更清楚嘛。我今天还在问那个谁，你对英伟达怎么看，它现在有没有可能被替代等等，我确实不懂。

% 来源：EPUB part256.xhtml
\BookMinorHeading{最终你会成为本该成为的人}

\BookQuestion{听众：这个时代机遇不等人，我是继续追求学业，还是为了抓住机遇而放弃学业，去走创业这条路？}

创业不知道，这取决于你。我认识一个人，也算是球友，后来在LSI做了。他在斯坦福大学读博士的时候，曾经有一个企业的老总跟他聊了两个小时，想把他挖走。但他是台湾人，他觉得我们中国人一定要先完成学业。跟他聊两个小时的那个人叫比尔·盖茨，结果他没有去。他一辈子都在后悔这件事，我猜的啊，我不知道。所以这个完全是你的决定。但是如果他碰见的是“贾布斯”，他又去了，那就是另外一个故事了，学业也没完成。

所以你要自己去判断什么是对的事情。要想长远，想10年、20年，而不是想着眼前这点利益。拿学位也是对的，但是，比尔·盖茨这个人跟你聊了两个小时，你都没有被打动，确实是愚蠢的，这个毫无疑问。就完全不知道对方是个什么样的人嘛，脑子只想着那一点点的小利益，要拿一个博士学位啊，不然对不起家长，可是硕士都已经有了嘛。当然那个时候微软还不是很厉害，他如果去了微软，可能就是微软前十号的员工。那他得多厉害，就是比现在的CEO，比李开复这些人都要早了。

\BookQuestion{听众：所以，人生还是有很多的机遇和运气。}

反正最终你会成为你本该成为的人。

% 来源：EPUB part257.xhtml
\BookMinorHeading{黄峥不是一个趋利的人}

\BookQuestion{主持人：拼多多和步步高、OPPO、vivo都把本分当作企业文化，但是在企业经营上有很大的不同，请问应该如何理解这种差异？}

我们公司我很清楚，拼多多我确实不太知道他们具体的东西。很多人跟我讲，我说人家的用户越来越多，那么多人喜欢用它，你非说它不对，到底是啥意思？我知道黄峥最早开始创业的时候，比方他做拼好货时，他是非常关心农业的，他就是想把农产品发出去。很多事情，包括他现在做的很多基础建设，真的都是跟那个有关的。

做拼多多的时候，是他找我，最后我投了他。我事先其实并不知道他在做什么，就有一天通电话，我说你在干嘛，他说他在做这个，我说这不就是以前的拼好货嘛？他说对，但是它现在有点不一样了，做的东西多一点了，等等吧。然后他说他希望我能够投他们。我问，那它能挣钱吗？他说不知道，但是用户量成长非常快，很多人喜欢用。而且供应商也很开心，原来不好卖的东西都卖出去了。尤其是农产品，因为他们最早是做农产品出现的。你想，农民的橘子如果卖不动，就都烂掉了，所有的水果都是一样，但是通过他们的渠道，确实卖得很好。他让我加入，我说能赚钱吗？他说不知道，我听他讲了成长的速度以后，说，那我就把它当公益做吧，因为那么快的成长速度，说明你是在做一件好事。

最后的结果就是这样，我就说如果赚钱了，我就把它捐到我的基金里头，就当做公益了；如果不赚钱，但是帮了那么多人，那我也当做公益了，反正都是做公益，一样。然后我给他提了个条件，说你想让我投多少，你就再投多少。所以他并不是因为缺钱找我的，他就是为了把我拉进去，做他的投资人。这个故事在很早年，是上市之前的事情，跟上市没有关系。

我觉得他们做的事情，第一我不是很了解细节。有些东西其实我也跟他提过，我说，哎，我怎么觉得你这个东西不对。他说，嗯，这个东西确实不对，他会改。但是我没有太跟踪很多细节，我看网上有各种各样的说法。但是事实上他确实厉害，营业额、用户量等等一直在往上涨。这些那么不喜欢的人也没试过，很多人从来没用过，就觉得那就是坏。我很难理解这种事情。

\BookQuestion{听众：拼多多供应商受到各种刁难……我听说了，首先我不好评价啊，因为我不了解具体的情况。但是大部分供应商还是认真地在那里待着。这是第一。第二他们的竞争环境可能更恶劣一些。我们的供应商很少，相比于他们而言。很多供应商有意见，你告诉我有多少。如果是有一万个供应商有意见，而它可能有五百万个供应商。所以从这个角度可能能理解，但是具体怎么回事，我不知道，所以我没有办法回答你这个问题。}

我觉得黄峥不是一个趋利的人，很多人不理解这一点，我非常理解。他不是为了挣钱，他也不会蒙你。前段时间拼多多股价狂掉的时候，很多人说，这有没有可能是个骗子。我说别的我不知道，这一点我知道，他一定不是。但是你喜不喜欢，那是你的问题，是你的选择。

% 来源：EPUB part258.xhtml
\BookMinorHeading{大学生迷茫很正常}

\BookQuestion{听众：我是大一的学生，有时候会迷茫，我大学这几年都要做什么？然后我读到黄峥的文章，他说在追求第一和努力做一个好学生上，浪费了太多的时间，损失了很多逆反捣蛋和纯粹享受青春的时光。60分万岁是个好哲学……但另一方面，如果我想加入一家我喜欢的比较好的公司，它肯定也会对我提出一些硬性的要求。这两者之间应该怎样平衡？}

大学生会迷茫很正常。我大概到大三的时候，才有一天突然发现，我准备考大学的过程中很开心、很充实。上了大学以后就很迷茫，尤其是我学的专业还不是我真的未来想干的事情。所以到大三的下半个学期，我才突然意识到，我的乐趣其实是来自过程。同时我也在大三的时候发现“要做对的事情，把事情做对”这一句话。反正，迷茫是正常的啊。

你说你喜欢的公司，对方怎么要求你？这个啥意思，我还没有完全理解，你能再重复一下吗？

\BookQuestion{听众：比方说，现在一个学生去应聘，公司可能对你的绩点等各方面有一个比较规范的要求。}

这个我不知道，因为每个公司的要求可能不一样。你喜欢那个公司，你不是得知道他们喜欢你什么吗？我怎么会知道呢？所以是需要你自己去了解的。我确实见过很多人去找工作，说我喜欢你们公司。我说，那你喜欢啥呀？他说不出来，不就是想找份工作嘛。你喜欢那家公司，一定有你喜欢的理由的。

跟你去做什么工作也有关系，比方说，因为你学的是软件工程专业，那你该学AI了，AI将来就替代了你的工作。吓唬你的，说着玩。我不知道，我真的是觉得，你自己要想，五年、十年、二十年以后，如果你很大的专长是写代码，那你可能真的需要找点别的事做。如果你足够聪明，你肯定是可以找得到的。我见过好多人改行，改到边儿都找不着的，但是人家一样可以做得很成功。你看我们班的龚剑萍，到日本做的事业就跟原来学的一点关系都没有，结果人家成了世界级的专家。

% 来源：EPUB part259.xhtml
\BookMinorHeading{如何培养能力圈}

\BookQuestion{听众：想做职业投资的新人，如何培养能力圈？}

你上大学不就是在培养你的能力圈嘛，40岁以前，你一定要想办法建立你的能力圈，但是在这以后，我觉得可能就会难。

\BookQuestion{听众：先在实业企业里面实践更有帮助，还是直接做投资？}

能在企业里面，当然是最好了。但是巴菲特也没有做过实业。
